% ============================================================================
% Część 1 przewodnika: produkt (wprowadzenie, problem, architektura, rekord,
% agent prowadzący wywiad, złożenie zgłoszenia). Plik jest wczytywany przez
% \input z przewodnik.pl.tex -- zawiera wyłącznie treść od \section.
% Źródła: README.md, docs/release/RELEASE-GATE.md, docs/architecture/*.md,
% docs/OPEN-PROBLEMS.md, docs/adr/0007. Rysunki: docs/diagrams/p1-*.mmd i system.mmd.
% ============================================================================

\section{Wprowadzenie}

\subsection{Czym jest ten produkt}

\textbf{Exit Interview Agent} to otwartoźródłowy agent AI, który prowadzi
ustrukturyzowane wywiady pożegnalne (ang.\ \emph{exit interview}) z byłymi
pracownikami. Z każdej rozmowy powstaje \emph{rekord}: krótki, ustrukturyzowany
dokument z ocenami sześciu tematów i kilkoma dosłownymi cytatami z wypowiedzi
rozmówcy. Rekordy wielu osób można następnie zsumować w porównywalne
\emph{sygnały o pracodawcy} --- ale tylko wtedy, gdy w każdej pokazywanej
komórce jest ich dość dużo i zawsze razem z miarą niepewności. Trzecim
składnikiem jest zestaw testów jakości \emph{samego wywiadu} (nie pracodawcy
i nie człowieka), oparty na śladach OpenTelemetry; jego opis należy do
dalszych części przewodnika.

Projekt sam o sobie mówi, że jest \emph{narzędziem i pracą portfolio},
a nie firmą ani publiczną platformą z opiniami (\repopath{README.md},
\repopath{docs/architecture/PROJECT-BRIEF.md}, §1). Jedna zasada wynika z tego
zdania i przewija się przez cały przewodnik: \finding{projekt nie hostuje żadnego
modelu językowego --- moc obliczeniową przynosi użytkownik} (własny klucz API
albo model lokalny).

\subsection{Dla kogo}

\begin{tabularx}{\linewidth}{@{}lX@{}}
\toprule
Rola & Co z tego ma i co robi \\
\midrule
Były pracownik &
Przechodzi rozmowę prowadzoną przez AI, która na początku mówi, że jest AI,
wyjaśnia, co zostanie zapisane, i prosi o jasną zgodę. Może przerwać w każdej
chwili --- przerwanie nie zostawia nic. Dostaje kod paragonu, którym później
może usunąć swój rekord. \\
Pracodawca (i czytelnik wyników) &
Nie dostaje pojedynczych opinii --- te nie są nigdy publikowane. Widzi tylko
zagregowane sygnały: z miarą niepewności, z podaniem pokrycia obok każdej oceny,
bez jednej łącznej punktacji i bez rankingu. \\
Deweloper / operator &
Uruchamia system lokalnie (jedno polecenie), podpina model, konfiguruje
tożsamość i księgę zgłoszeń, rozwija agenta i uruchamia testy. \\
\bottomrule
\end{tabularx}

\subsection{Mapa przewodnika}

Przewodnik ma sześć części. Ta pierwsza odpowiada na pytanie \emph{co i dlaczego}:
problem, zasady projektowe, architekturę, rekord, agenta i drogę rekordu do
serwera. Pozostałe części (obsługa, prywatność, jakość, ćwiczenia) opisują, jak z tego
korzystać, co chroni dane i jak mierzy się jakość wywiadu. Część szósta opisuje
wersję webową: drogę osoby, która rozmawia w przeglądarce, to, co dzieje się z danymi,
bramki kosztów oraz to, czego jeszcze nie wykonano (przegląd prawny, prawdziwy model, wdrożenie). Źródłem prawdy dla
każdej liczby i każdego twierdzenia są pliki w repozytorium (\repopath{docs/},
\repopath{README.md}, kod w \repopath{src/} i \repopath{web/}); gdzie czegoś nie
zmierzono, piszemy to wprost.

\subsection{Status projektu: co jest, a czego jeszcze nie ma}

README stosuje rygorystyczną konwencję: słowo \emph{Zaimplementowane} oznacza
kod na gałęzi \texttt{main}, wszystko inne jest planem przypisanym do zadania
z backlogu. W chwili pisania wszystkie zadania T0--T11 są oznaczone jako
zaimplementowane, a T12 (przegląd bezpieczeństwa, bramka wydania, opis wyników)
jako planowane. Poniższa tabela streszcza stan README.

\begin{longtable}{@{}p{0.31\linewidth}p{0.64\linewidth}@{}}
\toprule
Obszar & Stan \\
\midrule
\endhead
Szkielet (Aspire AppHost, wspólne jądro, jedna usługa z własną bazą, portal Next.js z BFF, kontenery, CI, skan sekretów) & zaimplementowane (T0) \\
Tożsamość: \texttt{authservice} jako przypięty obraz, walidacja JWT tylko RS256 & zaimplementowane (T0) \\
Rekord v1, biblioteka walidacji, deterministyczny detektor PII & zaimplementowane (T1) \\
Dwa schematy JWT (web i MCP), rotacja odświeżania w BFF, zgody, usuwanie konta & zaimplementowane (T2); nie uruchomiono: pełnego przepływu z Claude i obrazu \texttt{v0.3.4} \\
% zweryfikowano: README.md:26 (wiersz „Documentation foundation”, T3: zaimplementowane, tylko dokumenty)
Dokumentacja fundamentu (prywatność, zagrożenia, prawo, otwarte problemy, metodologia) & zaimplementowane (T3); tylko dokumenty, opisują projekt \\
Rdzeń agenta, skryptowy model testowy, osiem symulowanych person, demo offline & zaimplementowane (T4); model testowy to element testów, nie wzorzec jakości \\
Zgłoszenia po stronie serwera: walidacja, księga, kod paragonu, bilety, retencja & zaimplementowane (T5); weryfikator zatrudnienia to atrapa, niczego nie weryfikuje (OP-1) \\
Dostawcy modeli (Anthropic, zgodne z OpenAI, Ollama), interaktywne \texttt{interview} & zaimplementowane (T6); \textbf{nie wywołano żadnego prawdziwego dostawcy} \\
Zestaw testów jakości (27 scenariuszy, asercje warstwy 1, sędzia warstwy 2, baseline, bramka CI) & zaimplementowane (T7); tylko profil z modelem testowym, \textbf{żaden prawdziwy model nie był oceniany} \\
Serwer MCP (tryb A) & zaimplementowane (T8); \textbf{nie uruchomiono z prawdziwym klientem Claude} (OP-18) \\
Panel web & zaimplementowane (T9); testowany względem atrapy backendu, nie prawdziwych usług (OP-17); tylko angielski; bez ręcznego przeglądu dostępności (OP-16) \\
Sygnały (moduł i strony web) & zaimplementowane (T10, T10b); strony testowane względem atrapy zgodnej z kontraktem (OP-28); zrozumiałość dla czytelników nie była testowana (OP-26) \\
Zgłoszenie z CLI (\texttt{submit}, \texttt{delete-receipt}) & zaimplementowane (T11); nie uruchomiono względem wdrożenia ani przez prawdziwe TLS (OP-23 do OP-25) \\
Przegląd bezpieczeństwa, bramka wydania, opis wyników & planowane (T12) --- patrz niżej \\
\bottomrule
\end{longtable}

\uwaga{Nic nie jest wdrożone i w tej fazie nic nie będzie. Pliki Fly.io
i workflowy \texttt{flyio*.yml} są wygenerowane, ale nigdy nie uruchomione;
repozytorium jest prywatne, a upublicznienie jest decyzją właściciela.
\textbf{Projekt nie przetwarza danych żadnej prawdziwej osoby} i nie oferuje
prawdziwych wywiadów --- dopóki nie istnieje polityka prywatności, podstawa
prawna, działające usuwanie i przegląd prawnika (README, sekcja o celach
niebędących celami).}

Dokument \repopath{docs/release/RELEASE-GATE.md} (uruchomiony 2026-10-05) wydaje
werdykt \textbf{NIE GOTOWE do publicznego wydania}. Pozycje bramki mają wynik
PASS, FAIL, NOT RUN lub OWNER DECISION. Z tabeli bramki wynika:

\begin{itemize}[itemsep=2pt]
  \item \textbf{PASS}: brak podatnych pakietów NuGet; obecność plików LICENSE,
        SECURITY.md i innych; poprawność odnośników w dokumentacji; plakietki
        tylko dla rzeczy, które istnieją.
  \item \textbf{FAIL}: 16 commitów na \texttt{main} zawiera nazwę modelu
        w komunikacie lub stopce (wbrew regule projektu --- decyzja właściciela:
        zostawić albo przepisać historię); a także proces --- trzy dokumentacyjne
        commity trafiły bezpośrednio na \texttt{main} bez PR i CI.
  \item \textbf{OWNER DECISION}: jedna wysoka podatność npm (\texttt{braces}),
        obecna tylko w zależnościach deweloperskich; widoczność repozytorium
        pozostaje niezmieniona.
  \item \textbf{NOT RUN}: skan sekretów całej historii prawdziwym skanerem,
        audyt licencji zależności, weryfikacja źródeł prawnych (proxy odmówiło
        dostępu do 9 źródeł), a także sprawdzenia ``w świecie rzeczywistym'':
        żywy konektor Claude, logowanie dwuskładnikowe z prawdziwym
        \texttt{authservice}, prawdziwe TLS, Windows i macOS, przegląd
        z czytnikiem ekranu, ocena prawdziwego modelu.
\end{itemize}

\finding{Dokument bramki wyraźnie zastrzega, że to nie jest dowód na obecność
sekretów w repozytorium --- to wykaz tego, co sprawdzono i czego nie.}

\section{Problem i zasady projektowe}

\subsection{Problem, który produkt próbuje rozwiązać}

Wywiad pożegnalny prowadzony przez pracodawcę ma wrodzoną wadę: osoba, która
odchodzi, rzadko mówi wprost to, co myśli o przełożonym, i trudno to potem
porównać między firmami. Dokumentacja projektu nie opiera się jednak na
obietnicy ``szczerości'' --- opisuje cel bardziej skromnie: ustrukturyzowany
wywiad, który daje \emph{porównywalny} rekord, oraz agregaty, które nie
pozwalają zidentyfikować pojedynczego autora. Dlatego ciężar projektu leży
nie w samej rozmowie z modelem, lecz w tym, \emph{co się z nią dzieje dalej}:
co jest zapisywane, czego celowo nie zapisuje się wcale i kto może coś
połączyć z czym.

\subsection{Zasady projektowe wynikające z briefu}

Dokument \repopath{docs/architecture/PROJECT-BRIEF.md} jest wiążący dla każdego,
kto pracuje nad repozytorium. Zasady, które kształtują produkt, to:

\begin{enumerate}[itemsep=3pt]
  \item \textbf{Konto nie ma związku z pracodawcą.} Logowanie do portalu
        identyfikuje tylko konto. Pracodawca jest polem zgłaszanego rekordu;
        nigdy nie modeluje się zdania ``to konto należy do pracodawcy X''.
  \item \textbf{Rekord nie ma identyfikatora osoby.} Brak identyfikatora
        użytkownika, adresu e-mail, IP, nazwiska i znacznika czasu (więcej
        w rozdziale o rekordzie).
  \item \textbf{Jeden wpis na pracodawcę na konto.} Wymusza to osobna
        \emph{księga zgłoszeń}, w której nie ma treści ani identyfikatora rekordu,
        tylko zaszyfrowany kluczem skrót (HMAC) pary: konto i pracodawca.
  \item \textbf{Usuwanie przez kod paragonu.} Po zgłoszeniu użytkownik dostaje
        losowy kod; serwer trzyma tylko jego skrót; okazanie kodu usuwa rekord
        bez łączenia go z kontem.
  \item \textbf{Agregaty tylko od progu K.} Domyślnie K wynosi 5 i dotyczy
        \emph{każdej pokazywanej komórki}, a nie tylko całego pracodawcy;
        publikacja jest wsadowa; zawsze podaje się niepewność; nie ma jednej
        łącznej oceny pracodawcy ani rankingu.
  \item \textbf{Rekordy są danymi osobowymi.} Nigdy nie nazywa się ich
        ``anonimowymi'' ani ``zanonimizowanymi'' (ADR-0018); agregaty można
        nazywać ``zagregowanymi''.
  \item \textbf{Żadnych danych osobowych ani treści rozmów w logach, śladach
        i zdarzeniach audytu.}
  \item \textbf{Każdy rekord od klienta jest niezaufanym wejściem}: walidacja
        schematu, wykrywanie PII, limity szybkości i rozmiaru.
  \item \textbf{Zasady agenta:} brak pytań sugerujących; prośba o konkretny
        przykład, gdy twierdzenie jest mgliste; nigdy nie pytać o nazwiska ani
        ich nie zapisywać (wykryć i zamaskować); respektować wycofanie zgody
        i przerwać na żądanie; na początku ujawnić, że to AI; w schemacie
        ekstraktora nie ma pola na emocje, sentyment ani afekt.
\end{enumerate}

\subsection{Czego projekt jawnie nie robi (cele negatywne)}

Brief zawiera listę ``non-goals'', które są równie ważne jak cele:

\begin{itemize}[itemsep=2pt]
  \item żadnego blockchaina, tokenów ani DAO (niezmienność kłóci się z prawem
        do usunięcia danych w RODO);
  \item żadnej publikacji pojedynczych opinii --- tylko agregaty powyżej minimum;
  \item żadnych prawdziwych wywiadów z prawdziwymi ludźmi w tej fazie: wyłącznie
        symulowane persony;
  \item żadnej prawdziwej weryfikacji zatrudnienia: jest interfejs
        \texttt{IEmploymentVerifier} i atrapa, a prawdziwa weryfikacja jest
        opisana jako problem otwarty;
  \item brak obsługi tokenów subskrypcji Claude (Free/Pro/Max) ani GitHub
        Copilot jako backendu modelu. Dla Claude powodem jest dokumentacja
        Anthropic (czytana 2026-10-05), według której podmioty trzecie nie
        mogą oferować logowania Claude.ai ani kierować zapytań przez
        poświadczenia planów subskrypcyjnych. Dla Copilota powodem jest to, że
        warunków \emph{nie udało się zweryfikować} --- to nie jest ustalenie,
        że są zabronione;
  \item żadnego wdrożenia i żadnego upubliczniania repozytorium przez
        sesje robocze.
\end{itemize}

Obsługiwany dostęp do modelu to: klucz API Anthropic, dowolny punkt końcowy
zgodny z OpenAI i modele lokalne przez Ollama. Dla klucza Anthropic README
zauważa, że zastosowania związane z zatrudnieniem podlegają wymaganiom
Anthropic dla zastosowań wysokiego ryzyka; warunków OpenAI \emph{nie
udało się odczytać} (są oznaczone jako niezweryfikowane).

\subsection{Problemy otwarte --- co zostaje nierozwiązane}

Plik \repopath{docs/OPEN-PROBLEMS.md} wylicza 29 problemów (OP-1 do OP-29).
Przy każdym podaje: dlaczego ma znaczenie, co robi się teraz i co by go zamknęło.
Zasada jest taka, że wymieniony problem to świadoma decyzja, a niewymieniony
to dryf. Tabela poniżej pokazuje te, które najbardziej dotyczą tej części.

\begin{longtable}{@{}p{0.09\linewidth}p{0.62\linewidth}p{0.2\linewidth}@{}}
\toprule
Nr & Problem & Waga \\
\midrule
\endhead
OP-1 & Brak prawdziwej weryfikacji zatrudnienia --- system nie odróżni
       prawdziwego byłego pracownika od skryptu, konkurenta czy samego pracodawcy;
       każdy agregat jest ``deklarowany przez konta'' & wysoka \\
OP-2 & Brak rejestru pracodawców (wolny tekst rozbijałby jednego pracodawcę
       na wielu, co obchodzi próg K) & wysoka \\
OP-3 & Granulacja pasm a małe grupy & wysoka \\
OP-4 & Próba skrzywiona w stronę użytkowników technicznych --- do uruchomienia
       trzeba hosta MCP lub CLI z kluczem & średnia \\
OP-5 & Stronniczość modeli (inne pytania lub maskowanie nazwisk zależnie od
       języka i pochodzenia) & średnia \\
OP-6 & Wielojęzyczność --- stan to ``angielski testowany, polski nietestowany'' & średnia \\
OP-7 & Logowanie CLI bez klientów publicznych w \texttt{authservice} --- stąd
       bilety & średnia \\
OP-12 & Kod paragonu: dostęp bez listy rekordów; zgubiony kod oznacza brak
        możliwości usunięcia & średnia \\
OP-13 & Korelacja na poziomie składowania między księgą a rekordami & średnia \\
OP-18 & Tryb A: wierność hosta, tekst otwierający i niezweryfikowane
        uruchomienie z Claude & wysoka \\
OP-19 & K jest konwencją, a małe partie ujawniają małe różnice & wysoka \\
\bottomrule
\end{longtable}

\finding{Najpoważniejsze ograniczenie jest wpisane w sam produkt: operator,
który ma bazę, klucz podpisujący i żywy ruch, może skorelować konta z rekordami,
a małe grupy mogą być rozpoznawalne.} Dokumentacja prywatności i model zagrożeń
(część 3 przewodnika) opisują to szczegółowo.

\section{Architektura}

\subsection{Kształt systemu}

Dokument \repopath{docs/architecture/00-ARCHITECTURE.md} opisuje system jako:
jeden korzeń kompozycji (AppHost), jedno wspólne jądro, jedną usługę z własną
bazą, jedną powierzchnię produktu i jedną zużywaną usługę tożsamości.
Odniesieniem są zasady P1--P15 z repozytorium \texttt{architecture-}\allowbreak\texttt{standards};
każde odstępstwo to ADR i wiersz w rejestrze odstępstw z datą i powodem.

\begin{figure}[H]
\centering
\includegraphics[width=0.8\linewidth,height=0.5\textheight,keepaspectratio]{\dgm{system}}
\caption{Architektura systemu (\texttt{docs/diagrams/system.mmd}): przeglądarka
rozmawia tylko z portalem (bramka brzegowa i BFF); BFF wywołuje
\texttt{authservice} oraz \texttt{interview-service} z tokenem z ciasteczka
HttpOnly; klient MCP (Claude) uwierzytelnia się w \texttt{authservice}, a potem
wywołuje \texttt{/mcp}. AppHost jedynie komponuje całość (tylko rozwój).}
\end{figure}

\subsection{Projekty w \texttt{src/}}

\begin{longtable}{@{}p{0.38\linewidth}p{0.57\linewidth}@{}}
\toprule
Projekt & Rola \\
\midrule
\endhead
\repopath{AppHost} & korzeń kompozycji (P1), tylko do rozwoju; jedno polecenie
  uruchamia Postgres, \texttt{authservice}, \texttt{interview-service} i web \\
\repopath{ServiceDefaults} & wspólne jądro (P2): telemetria, zdrowie, wykrywanie
  usług, odporność, walidacja JWT, CORS, ograniczanie szybkości --- wyłącznie
  ``instalacja'', bez logiki domenowej; rozmiar pilnuje skrypt
  \texttt{check-kernel-size.sh} \\
\repopath{Contracts} & obiekty DTO i stabilne kody odrzuceń przekraczające
  granicę usług \\
\repopath{InterviewService} & jedyna usługa; właściciel bazy \texttt{interviewdb};
  zgłoszenia, księga, paragony, bilety, retencja; gości moduł Signals i punkt
  MCP \\
\repopath{Signals} & moduł sygnałów: reguły ujawniania, statystyka, wsadowy
  publikator, własny schemat bazy; zna tylko jądro \\
\repopath{Records} & rekord wywiadu: niezmienny model, walidacja schematu,
  postać kanoniczna, wierność cytatów \\
\repopath{Privacy} & deterministyczny detektor i maskownik PII \\
\repopath{Agent} & rdzeń agenta: protokół, maszyna stanów, role, strażnik PII,
  krok cytatów, model testowy \\
\repopath{Personas} & symulowani rozmówcy sterowani danymi i ziarnem losowości \\
\repopath{Providers} & dostawcy modeli za \texttt{IChatClient} \\
\repopath{Cli} & program \texttt{exit-interview}: demo, \texttt{interview},
  \texttt{providers}, \texttt{submit}, \texttt{delete-receipt} \\
\repopath{Eval} & zestaw ocen jakości wywiadu i narzędzie \texttt{exit-interview-eval} \\
\bottomrule
\end{longtable}

Poza \texttt{src/} istnieją: \texttt{web/app} (portal Next.js z BFF),
\texttt{schemas/} (publikowane, wersjonowane schematy: rekord, wyjście
ekstraktora, persona), \texttt{tests/} (projekty xUnit odpowiadające źródłom
oraz scenariusze Playwright w \texttt{tests/e2e}) i \texttt{flyio/}
(wygenerowana, niewdrożona topologia). Liczby zależności bezpośrednich
projektów (np.\ 0 dla \texttt{Contracts} i \texttt{Privacy}, 12 dla jądra)
pochodzą z tabeli w \texttt{00-ARCHITECTURE.md} i są odtwarzalne poleceniem
\texttt{grep -c PackageReference} na plikach projektów.

\subsection{Usługa tożsamości \texttt{authservice}}

Projekt \emph{nie} implementuje logowania. Korzysta z osobnej usługi
\texttt{authservice} jako przypiętego obrazu kontenera (z własną bazą
\texttt{authdb} i własnym kluczem podpisującym; kod nigdy nie jest kopiowany).
To \texttt{authservice} utrzymuje jedyny klucz podpisujący (RS256), a usługi
walidują tokeny wyłącznie tym algorytmem, na podstawie zestawu kluczy JWKS.
Rejestruje ona tylko klientów \emph{poufnych} OAuth (sekret co najmniej
32 bajty, bez klienta publicznego, bez przepływu urządzenia, bez dynamicznej
rejestracji). Ma to dwie konsekwencje, które kształtują architekturę:

\begin{itemize}[itemsep=2pt]
  \item \textbf{Dwa schematy JWT} w \texttt{interview-service}: zwykły
        (web/BFF) oraz schemat MCP, w którym token jest związany z odbiorcą
        (\texttt{aud}) zasobu MCP, a zakres (\texttt{interview:submit})
        egzekwuje sama usługa. Tokeny jednego schematu są odrzucane przez drugi
        (testy negatywne macierzy tokenów).
  \item \textbf{CLI nie używa OAuth.} Opublikowany plik binarny nie utrzyma
        sekretu, więc zamiast tego portal mintuje jednorazowy bilet (rozdział~6).
\end{itemize}

Gdy obraz \texttt{authservice} nie jest dostępny, stos można uruchomić bez
tożsamości (\texttt{Identity\_\_}\allowbreak\texttt{Enabled=false}); chronione punkty końcowe
odpowiadają wtedy 401 i mówią o tym wprost. Zasada P8 (opcjonalne zależności
degradują się) jest realizowana przez \texttt{IntegrationStatus}: punkt
\texttt{/health} wylicza, co działa (tożsamość, uwierzytelnianie MCP, baza,
eksport telemetrii, klucz księgi, weryfikator zatrudnienia, sygnały), a ten sam
spis wypisuje baner startowy.

\subsection{Portal web i BFF}

\texttt{web/app} jest zarazem powierzchnią produktu i BFF (\emph{backend for
frontend}). Przeglądarka rozmawia wyłącznie z tym samym źródłem; tokeny
trzymane są w ciasteczkach HttpOnly (SameSite=Strict), a nie w skryptach strony.
Najważniejsze mechanizmy (opis w tabeli ``The web panel'' w dokumencie
architektury):

\begin{itemize}[itemsep=2pt]
  \item \textbf{Bramka brzegowa} (\texttt{proxy.ts}): weryfikuje podpis,
        wystawcę i odbiorcę tokenu, ma jawną listę ścieżek publicznych, wymusza
        krok zgody, ustawia CSP z nonce na każde żądanie, \texttt{no-store}
        i kontrolę tego samego źródła dla zmian stanu.
  \item \textbf{Proxy z jedną ścieżką} (\texttt{app/api/proxy}): mapuje tylko
        ścieżki \texttt{v1/*} i wstrzykuje token z ciasteczka.
  \item \textbf{Osobna anonimowa ścieżka usuwania przez paragon}, bez konta
        i bez tokenu --- bo ten punkt końcowy jest anonimowy z założenia (ADR-0049).
  \item \textbf{Brak listy ``moje zgłoszenia''} --- celowo i na stałe: rekordy
        nie są powiązane z kontami, więc takiej listy nie dałoby się zbudować
        bez zerwania zasady prywatności.
\end{itemize}

Portal nie uruchamia żadnego modelu. Ekrany: strona główna, logowanie z drugim
składnikiem, rejestracja i weryfikacja e-mail, zgoda, podłączenie klienta AI,
bilet CLI (pokazywany raz), usuwanie przez kod paragonu, konto (eksport danych,
usunięcie), strona prywatności oraz strony sygnałów. Portal jest po angielsku
i ma bramkę dostępności axe-core; ręcznego przeglądu dostępności nie
wykonano (OP-16).

\subsection{Trzy tryby użycia}

Brief definiuje trzy tryby. Wspólna cecha: \emph{serwer nigdy nie widzi
zapisu rozmowy, tylko rekord.} Różnią się tym, kto rozmowę prowadzi i kto ją
widzi.

\begin{figure}[H]
\centering
\includegraphics[width=0.85\linewidth,height=0.45\textheight,keepaspectratio]{\dgm{p1-tryby}}
\caption{Trzy tryby: kto prowadzi rozmowę i co trafia do serwera. W trybie A
cała rozmowa widoczna jest dla Claude i jego dostawcy; w trybie B --- dla
dostawcy modelu, który wybrałeś (albo dla nikogo, przy modelu lokalnym).}
\end{figure}

\begin{longtable}{@{}p{0.17\linewidth}p{0.38\linewidth}p{0.4\linewidth}@{}}
\toprule
Tryb & Jak działa & Kto widzi rozmowę \\
\midrule
\endhead
\textbf{A: MCP z Claude} &
Użytkownik łączy Claude (własny klient Anthropic) z serwerem MCP.
Model-gospodarz prowadzi wywiad według \emph{promptu}, który serwer publikuje,
a rekord wysyła narzędziem. Obsługiwany jest wyłącznie Claude (redirect URI
\texttt{claude.ai}); inne hosty to praca przyszła (OP-9). &
Claude i jego dostawca widzą całą rozmowę, serwer --- tylko rekord. README
nazywa ten tryb najsłabszym z trzech pod względem prywatności i wierności:
serwer nie obserwuje, czy host ujawnił AI, uzyskał zgodę, zadawał neutralne
pytania. \\
\textbf{B: CLI} &
Program \texttt{exit-interview} prowadzi cały wielorolowy wywiad lokalnie
(prowadzący, dopytujący, ekstraktor rekordu, strażnik PII) z kluczem API
użytkownika lub modelem lokalnym. Zapis rozmowy zostaje w pamięci na maszynie. &
Dostawca modelu, który wybrałeś, widzi całą rozmowę; przy Ollama na własnej
maszynie program niczego nie wysyła. Do serwera trafia tylko rekord
i bilet --- i tylko jeśli zdecydujesz się uruchomić \texttt{submit}. \\
\textbf{C: portal} &
Konto, zgody, bilety CLI, usuwanie przez paragon, prywatność oraz strony
sygnałów. Nie uruchamia żadnego modelu. & Nie prowadzi rozmów. \\
\bottomrule
\end{longtable}

\uwaga{Tryb offline (\texttt{demo}) uruchamia całą rozmowę z symulowaną osobą
i deterministycznym modelem testowym: bez sieci, bez kluczy, bez wysyłania.
Z tym samym ziarnem wynik jest identyczny bajt w bajt (sprawdzono: dwa
uruchomienia \texttt{demo --persona talkative --seed 1} dały identyczne
112 linii wyjścia, RESULTS.md §1).}

\section{Rekord wywiadu}

Rekord jest umową całego systemu: produkują go CLI i ścieżka MCP, waliduje
i przechowuje usługa przyjmowania, agreguje moduł sygnałów, a zestaw ocen
sprawdza jego wierność. Poniżej streszczenie dokumentu
\texttt{docs/architecture/}\allowbreak\texttt{record-schema.md} i ADR-0007. Schemat
(JSON Schema 2020-12, plik \texttt{record.v1.schema.json} w katalogu \texttt{schemas/})
oraz biblioteka \texttt{Records} są źródłem prawdy.

\finding{Rekord jest pseudonimowy, nie anonimowy (ADR-0018).} Zawiera
ustrukturyzowane oceny sześciu tematów, dosłowne cytaty uzasadniające je
i grube pasma kontekstu; nie zawiera niczego, co identyfikuje osobę lub łączy
z jej kontem.

\subsection{Pola rekordu}

\begin{longtable}{@{}p{0.3\linewidth}p{0.65\linewidth}@{}}
\toprule
Pole & Znaczenie \\
\midrule
\endhead
\texttt{schemaVersion} & ciąg \texttt{"1"}; wybiera umowę (tylko wersja główna) \\
\texttt{interviewId} & 32 małe znaki szesnastkowe = 128 losowych bitów z CSPRNG
  systemu; nie pochodzi z konta, pracodawcy, czasu ani treści \\
\texttt{employerRef} & nieprzezroczyste odniesienie do pracodawcy, 3--64 znaki,
  wzorzec \texttt{[a-z0-9]} z myślnikami; ścisły wzorzec wyklucza wolny tekst,
  nazwiska i adresy e-mail \\
\texttt{context.tenureBand} & pasmo stażu (wymagane) \\
\texttt{context.seniorityBand}, \texttt{context.functionBand} & pasma opcjonalne;
  gdy wstrzymane, są \emph{pominięte}, nigdy \texttt{null} \\
\texttt{topics.<temat>} & zawsze wszystkie sześć tematów:
  \texttt{onboarding}, \texttt{management}, \texttt{growth},
  \texttt{pay\_vs\_promises}, \texttt{culture}, \texttt{reason\_for\_leaving} \\
\texttt{status} & \texttt{no\_data} albo \texttt{covered}. \emph{Brak danych
  nie jest niską oceną} \\
\texttt{rating} & liczba całkowita 1--5 (1 bardzo negatywna, 5 bardzo
  pozytywna) lub \texttt{null} \\
\texttt{confidence} & \texttt{low}, \texttt{medium}, \texttt{high} lub
  \texttt{null}: trzy poziomy, nie liczba --- ułamek udawałby precyzję,
  której ekstrakcja nie ma \\
\texttt{quotes} & do 5 dosłownych cytatów po maks.\ 400 punktów kodowych \\
\texttt{piiMasked} & prawda, gdy cytaty pochodzą z zapisu, który przeszedł
  maskownik PII; przyjmowanie odrzuca \texttt{false} domyślnie \\
\texttt{interview.protocolVersion} & wersja protokołu wywiadu \\
\texttt{interview.language} & kod języka ISO 639 \\
\texttt{interview.aiDisclosed} & prawda, gdy prowadzący powiedział, że jest AI,
  przed rozpoczęciem; przyjmowanie odrzuca \texttt{false} domyślnie \\
\texttt{interview.durationBand}, \texttt{interview.turnBand} & pasma czasu trwania
  i liczby tur (po cztery wartości) \\
\bottomrule
\end{longtable}

Spójność tematu wymusza sam schemat: dla \texttt{no\_data} ocena i pewność są
\texttt{null}, a cytatów jest zero; dla \texttt{covered} jest od 1 do 5
cytatów, a ocena nie występuje bez cytatu wspierającego (inaczej kod błędu
\texttt{TOPIC\_INCONSISTENT}). Agregacja musi liczyć \texttt{no\_data} jako
brak, nigdy jako niski wynik --- dlatego to osobny stan.

\subsection{Pasma kontekstu}

Kontekst to dokładnie trzy pola, każde będące enumeracją --- klient nie może
wpisać prozy (ADR-0007).

\begin{tabularx}{\linewidth}{@{}lXl@{}}
\toprule
Pole & Wartości & Wymagane \\
\midrule
\texttt{tenureBand} & \texttt{lt\_6m}, \texttt{6m\_1y}, \texttt{1y\_3y}, \texttt{3y\_5y}, \texttt{5y\_10y}, \texttt{gt\_10y} & tak \\
\texttt{seniorityBand} & \texttt{junior}, \texttt{mid}, \texttt{senior}, \texttt{management} & nie \\
\texttt{functionBand} & \texttt{engineering}, \texttt{product\_design}, \texttt{sales\_marketing}, \texttt{operations\_support}, \texttt{corporate\_functions}, \texttt{other} & nie \\
\bottomrule
\end{tabularx}

\medskip
Uzasadnienie (ADR-0007): pełny iloczyn pasm to $6\times4\times6 = 144$ komórek.
Pasma są nieliczne i szerokie, tak by żadne nie znaczyło ``garstki ludzi'' u
średniego pracodawcy: \texttt{management} łączy kierowników zespołów, menedżerów
i dyrektorów (jednoosobowe pasmo ``dyrektor'' byłoby nazwiskiem),
\texttt{corporate\_functions} łączy finanse, prawo, HR i administrację, a
\texttt{other} istnieje, by nikt nie musiał wybierać błędnego pasma. Staż jest
jedynym wymaganym pasmem, bo jest najmniej identyfikujący, a najbardziej
przydatny przy interpretacji tematów. Otwarte górne pasmo uniemożliwia
singleton ``pracownik założyciel''. Metadane wywiadu też są w pasmach
(czas: od \texttt{lt\_10m} do \texttt{gt\_40m}; liczba tur: \texttt{lt\_10},
\texttt{10\_20}, \texttt{20\_40}, \texttt{gt\_40}), bo dokładny czas lub liczba
tur to odcisk palca sesji.

Limity są częścią kontraktu prywatności, a nie wyborem UX: 5 cytatów na temat,
400 punktów kodowych na cytat (``około trzech, czterech zdań'': mieści jeden
konkretny przykład, a nie długą identyfikującą anegdotę), co daje najwyżej 30
cytatów i 12\,000 punktów kodowych na rekord, oraz ładunek maks.\ 160~KiB.
Cytat nie może zawierać znaków sterujących, końców linii ani znaków
nadpisujących kierunek tekstu. ADR zastrzega: liczby te to osądy, nie
pomiary --- nic nie sprawdzono na prawdziwych wywiadach, bo ich nie ma.

\subsection{Czego w rekordzie celowo nie ma}

\begin{itemize}[itemsep=2pt]
  \item \textbf{Żadnego identyfikatora osoby}: identyfikatora użytkownika ani
        konta, e-maila, adresu IP, nazwiska, urządzenia ani sesji. Jedynym
        identyfikatorem jest losowy \texttt{interviewId}.
  \item \textbf{Żadnego znacznika czasu}: to klucz do łączenia z logami dostępu.
        Polityka prywatności dopuszcza najwyżej pasmo tygodnia ISO; schemat nie
        niesie żadnego, czyli jest bardziej rygorystyczny.
  \item \textbf{Żadnych emocji, sentymentu ani afektu.} Oceny dotyczą tematów,
        nigdy wniosków o uczuciach osoby; ta sama lista zakazanych słów jest
        sprawdzana testem architektury w schemacie rekordu i ekstraktora.
  \item \textbf{Bez dokładnych dat} (dołączenia, odejścia, wywiadu), wieku,
        płci i innych kategorii chronionych, lokalizacji, wielkości pracodawcy,
        identyfikatorów przełożonego lub zespołu, tytułu stanowiska i
        wynagrodzenia.
  \item \textbf{Brak pól dodatkowych}: każdy obiekt ma
        \texttt{additionalProperties: false}, więc nieznane pole nie może się
        przemycić. Błąd dla nieznanego klucza podaje ścieżkę rodzica, nigdy
        klucz (klucz mógłby być adresem e-mail).
\end{itemize}

Te zakazy są \emph{egzekwowane, nie obiecywane}: testy architektury oblewają
build, jeśli właściwość schematu lub składowa modelu przypomina identyfikator,
jeśli jakikolwiek obiekt schematu jest otwarty na dodatkowe pola albo jeśli
pole kontekstu nie jest enumeracją (ADR-0011).

\subsection{Cytaty: ryzyko resztkowe}

Cytaty to tekst wolny, więc to one są głównym wektorem reidentyfikacji.
Ograniczono je limitami, pobiera się je z zapisu zamaskowanego i sprawdza
wierność (poniżej). Maskowanie jest heurystyczne i ma zmierzone granice
(\texttt{docs/privacy/pii-detector.md}). Cytat nadal może zidentyfikować osobę
opisem (``jedyna kobieta w biurze w Gdańsku''); rekord tego nie naprawi i model
zagrożeń niesie to jako ryzyko resztkowe. Cytaty są też \emph{danymi
nieaktywnymi}: nic ich nie interpretuje, a fragment JSON, odwołanie \texttt{\$ref}
lub zdanie z wstrzyknięciem promptu wewnątrz cytatu jest tylko ciągiem znaków.
Konsument musi tę dyscyplinę zachować: kodować przy wyświetlaniu, nigdy nie
doklejać cytatu do zapytania ani promptu jako polecenia.

\subsection{Walidacja i wierność cytatów}

\texttt{RecordValidator.Validate} zwraca albo niezmienny \texttt{InterviewRecord},
albo listę błędów (kod, ścieżka). \textbf{Błędy nigdy nie zawierają
przesłanego tekstu.} Kolejność kontroli: rozmiar, składnia JSON i głębokość,
zduplikowane klucze, schemat, polityka \texttt{piiMasked}. Wybrane kody:
\texttt{PAYLOAD\_}\allowbreak\texttt{TOO\_LARGE}, \texttt{NOT\_JSON}, \texttt{DUPLICATE\_}\allowbreak\texttt{KEY},
\texttt{UNKNOWN\_}\allowbreak\texttt{FIELD}, \texttt{TOPIC\_}\allowbreak\texttt{INCONSISTENT}, \texttt{AI\_NOT\_}\allowbreak\texttt{DISCLOSED},
\texttt{PII\_NOT\_}\allowbreak\texttt{MASKED}, \texttt{VALIDATION\_}\allowbreak\texttt{TIMEOUT} (zabezpieczenie przed
atakiem na wyrażenia regularne). Kody są częścią umowy: zmiana lub ponowne
użycie tylko z nową wersją główną. Serializacja kanoniczna daje te same bajty
dla tego samego rekordu.

\texttt{QuoteVerifier} sprawdza, że każdy cytat jest dosłownym podciągiem
zapisu rozmowy po jednej normalizacji stosowanej do obu stron (ciągi białych
znaków zamieniane na jedną spację, końce przycięte). Wielkość liter,
interpunkcja i znaki diakrytyczne porównuje się porządkowo, więc parafraza,
zmiana wielkości liter czy zgubiony ogonek nie przejdą. \finding{Weryfikacja
działa tylko po stronie klienta i w ocenach: serwer nigdy nie ma zapisu rozmowy,
więc nie może sprawdzić cytatów.} Zmyślony cytat, który przechodzi schemat,
jest po stronie serwera nieodróżnialny od prawdziwego; wrogi klient może
przesłać cokolwiek (patrz model zagrożeń).

\subsection{Wersjonowanie}

\begin{itemize}[itemsep=2pt]
  \item \textbf{v1 jest niezmienna po wydaniu.} Po bramce wydania zmienia się
        tylko brzmienie opisów (\texttt{description}).
  \item \textbf{Zmiana addytywna = nowa wersja główna będąca nadzbiorem starej}:
        może dodać pola opcjonalne, poszerzyć enum, podnieść limit lub dodać
        opcjonalny temat; każdy ważny rekord v1 pozostaje ważny w v2, więc
        przechowywane rekordy nie wymagają migracji. Ponieważ walidatory v1
        odrzucają rekordy v2, wdrożenia aktualizują walidatory przed
        producentami.
  \item \textbf{Zmiana łamiąca} (usunięcie lub przemianowanie pola, zawężenie
        enum, obniżenie limitu, zmiana typu, uczynienie pola wymaganym,
        zmiana znaczenia \texttt{no\_data}) wymaga ADR, notatki migracyjnej
        i kolejności wdrożenia.
  \item \textbf{Żadna wersja schematu nie może zawierać identyfikatora osoby.}
  \item Biblioteka osadza schemat z pliku publikowanego, a test porównuje je
        bajt w bajt.
\end{itemize}

Otwarte kwestie rekordu: moduł sygnałów musi zastosować próg minimalnej liczby
do \emph{każdego} publikowanego przekroju, także w podziałach na pasma (schemat
tego nie wymusi); \texttt{employerRef} nie ma rejestru, więc dwa odniesienia
mogą nazywać jednego pracodawcę; nie zmierzono, czy wyjście prawdziwych modeli
wymaga normalizacji Unicode przed porównaniem cytatów.

\section{Agent prowadzący wywiad}

\subsection{Co to jest}

Rdzeń agenta (\texttt{src/}\allowbreak\texttt{ExitInterviewAgent.Agent}, opis w
\repopath{docs/architecture/interview-agent.md}) to rdzeń trybu B oraz
implementacja referencyjna protokołu, który tryb A udostępnia jako prompty
i zasoby. Prowadzi jeden ustrukturyzowany wywiad, maskuje to, co mówi
rozmówca, jeszcze przy odbiorze i --- jeśli zgoda trwała przez cały czas ---
zamienia zamaskowany zapis w zwalidowany \texttt{InterviewRecord}. Rdzeń
sam niczego nie wysyła (robi to \texttt{submit}, rozdział~6).

Naczelna zasada: \finding{program decyduje, co się dzieje; model tylko to
formułuje.} Wynika z niej cała budowa.

\begin{figure}[H]
\centering
\includegraphics[width=0.9\linewidth,height=0.45\textheight,keepaspectratio]{\dgm{p1-agent}}
\caption{Granice zaufania agenta (uproszczenie diagramu z
\texttt{interview-agent.md}). Surowa odpowiedź jest oczyszczana i maskowana
\emph{zanim} cokolwiek ją zapisze, przeanalizuje lub pokaże modelowi; maszyna
stanów, a nie model, wybiera każdy krok.}
\end{figure}

\subsection{Protokół}

Protokół jest wersjonowany (wersja \texttt{1.1}, ADR-0062) i przechowywany jako
dane w \texttt{interview-}\allowbreak\texttt{protocol.v1.json}, wbudowane w zestaw; te same dane
czytają adapter MCP i CLI. Zmiana jakiegokolwiek brzmienia lub limitu w tym
pliku jest zmianą zachowania protokołu i podnosi wersję (poboczną dla
brzmienia i limitów, główną dla zmiany tematów lub kolejności). Wersja trafia do
\texttt{interview.protocolVersion} każdego rekordu.

\textbf{Tura otwierająca to stały tekst, a nie wyjście modelu.} Mówi, że
prowadzący jest AI, co zostanie zapisane (krótki ustrukturyzowany rekord: ocena
i kilka własnych zdań rozmówcy na temat, losowy identyfikator, szerokie pasma;
bez nazwisk; bez powiązania z kontem; zamaskowane; sama rozmowa nie jest
zapisywana), że można przerwać w każdej chwili i że przerwanie niczego nie
zostawia, oraz prosi o jasne ``tak'' lub ``nie''. Ponieważ model nie może jej
napisać, nie może jej pominąć ani złagodzić. Flaga \texttt{aiDisclosed} jest
ustawiana na prawdę dopiero po tym, gdy runner dostarczył tę turę
\emph{i otrzymał odpowiedź}; jeśli nikt nie odpowie, przebieg kończy się jako
porzucony i rekord nie powstaje.

Sześć tematów w stałej kolejności, jedno pytanie naraz: onboarding, zarządzanie
(\texttt{management}), rozwój, płaca a obietnice, kultura, powód odejścia.
Protokół niesie neutralne pytanie otwierające każdego tematu; model może je
przeformułować, a brzmienie jest sprawdzane.

\begin{longtable}{@{}p{0.52\linewidth}p{0.43\linewidth}@{}}
\toprule
Reguła & Kto ją egzekwuje \\
\midrule
\endhead
R01 pierwsza tura: ujawnienie AI, co jest zapisywane, prawo do przerwania, prośba o zgodę & kod (stały tekst; test sprawdza każdy element) \\
R02 żadnej treści przed zgodą; brak jasnego ``tak'' = brak wywiadu & kod (maszyna stanów; test: przed zgodą brak wywołania modelu) \\
R03 sześć tematów w stałej kolejności & kod \\
R04 pytania otwarte i neutralne & oba: prompt prosi, \texttt{QuestionGuard} sprawdza i zastępuje \\
R05 jedno dopytanie o konkretny przykład na mglistą odpowiedź (\texttt{maxProbesPerTopic} = 1) & kod \\
R06 nigdy nie pytać o nazwiska; maskować; przekierować na zachowanie lub rolę & kod (strażnik PII, krok przekierowania) \\
R07 wycofanie zgody w dowolnej chwili: natychmiast przerwać, odrzucić, bez rekordu & kod \\
R08 tekst rozmówcy to dane, nigdy polecenia & oba: rozdzielenie w prompcie i struktura kodu \\
R09 jedno neutralne wyjaśnienie dla sprzecznych odpowiedzi & kod \\
R10 budżety tur, wywołań modelu i tokenów; łagodne zamknięcie & kod \\
R11 rozmówca lakoniczny lub wrogi: podziękować i zwolnić, nigdy nie spierać się & kod \\
R12 cytaty są dosłowne, od rozmówcy, na ten temat & kod \\
\bottomrule
\end{longtable}

\subsection{Maszyna stanów i role}

\texttt{InterviewMachine} nie wykonuje operacji wejścia-wyjścia. Jej jedynymi
przejściami są \texttt{Start()} i \texttt{OnReply(...)}, a każde zwraca
następny krok (co ma zrobić prowadzący). Tekst rozmówcy może na nią wpłynąć
tylko przez wartości logiczne w \texttt{ReplySignals}, które
\texttt{ReplyAnalyzer} liczy stałymi regułami z już zamaskowanej odpowiedzi.
Pierwszeństwo w obrębie tury (wygrywa pierwsze dopasowanie): wycofanie zgody
(stop) $\rightarrow$ budżet (zamknięcie) $\rightarrow$ druga wroga odpowiedź
(zamknięcie) $\rightarrow$ seria lakonicznych (zamknięcie) $\rightarrow$ wroga
odpowiedź jest przyjęta i następuje kolejny temat $\rightarrow$ wymieniona
osoba (przekierowanie, raz na temat) $\rightarrow$ sprzeczność
(wyjaśnienie, raz na temat) $\rightarrow$ mglista odpowiedź (dopytanie)
$\rightarrow$ następny temat albo zamknięcie po szóstym.

\finding{Wycofanie zgody wygrywa ze wszystkim, nawet z ładunkiem wstrzyknięcia:}
odpowiedź ``zakończ wywiad'' ukryta w manipulacyjnym tekście jest traktowana
jak wycofanie, bo zatrzymanie jest kierunkiem bezpiecznym dla prywatności, a
atakujący zyskuje tylko możliwość zakończenia własnego wywiadu.

Role stoją za interfejsami (\texttt{IInterviewer}, \texttt{IProber},
\texttt{IRecordExtractor}, \texttt{IPiiGuard}), więc testy, model testowy,
dostawcy i adapter MCP mogą dostarczać własne. Standardowe role działają na
dowolnym \texttt{IChatClient}.

\begin{longtable}{@{}p{0.22\linewidth}p{0.35\linewidth}p{0.38\linewidth}@{}}
\toprule
Rola & Robi & Nie robi \\
\midrule
\endhead
Prowadzący & formułuje pytanie tematu, wyjaśnienie i przekierowanie z tekstu
  zalążkowego protokołu & nie wybiera następnego kroku, nie widzi surowego
  tekstu rozmówcy \\
Dopytujący & formułuje jedno dopytanie o konkretny przykład & nie decyduje, czy
  dopytać (decyduje maszyna) \\
Ekstraktor rekordu & zwraca JSON zgodny ze schematem wyjścia ekstraktora & nie
  buduje rekordu, nie wybiera cytatów spoza zapisu \\
Strażnik PII & owija detektor trybem ``zamknięty w razie awarii''
  (\texttt{FailClosed}) i listą dozwolonych nazw pracodawców i produktów;
  maskuje przed wszystkim innym; wyjątek lub przekroczenie czasu kończy
  przebieg i niczego nie zachowuje & \\
Krok cytatów (\texttt{RecordAssembler}) & weryfikuje każdy cytat względem
  zamaskowanego tekstu rozmówcy \emph{z tego tematu}; odrzuconych nie wysyła & \\
\bottomrule
\end{longtable}

Każde pytanie sformułowane przez model przechodzi przez \texttt{QuestionGuard}:
niepuste, najwyżej 500 znaków, jedno pytanie (bez drugiego znaku zapytania
i bez współrzędnego drugiego zdania pytającego), bez fragmentów promptu,
bez PII, bez wzorców sugerujących, nacechowanych lub negatywno-polarnych,
bez zamkniętych początków, a dopytanie musi prosić o przykład. Odrzucone pytanie
zastępuje własne brzmienie protokołu (i jest liczone), więc niedbały lub
przejęty model sprawia, że wywiad jest bardziej bezbarwny, a nie inny.

\subsection{Od rozmowy do rekordu}

\begin{enumerate}[itemsep=2pt]
  \item Runner oczyszcza każdą odpowiedź (usuwa znaki sterujące i nadpisujące
        kierunek, zwija białe znaki, obcina do \texttt{maxReplyChars} = 2000)
        i \textbf{maskuje ją strażnikiem PII, zanim zostanie zapisana,
        przeanalizowana lub pokazana jakiemukolwiek modelowi}. Surowa odpowiedź
        nie jest zachowywana.
  \item Gdy dialog kończy się z nienaruszoną zgodą, \emph{ekstraktor} widzi
        wyłącznie zamaskowany zapis, w bloku danych, pod własnym promptem
        systemowym; najwyżej dwie próby, a ponowna próba niesie tylko \emph{kody}
        błędów.
  \item Wyjście jest parsowane względem schematu ekstraktora (zamknięte
        obiekty, bez pól na sentyment, emocję, afekt, nazwisko ani
        identyfikator).
  \item \texttt{RecordAssembler} buduje rekord: temat potrzebuje co najmniej
        \texttt{minWordsForCoverage} własnych słów rozmówcy (inaczej
        \texttt{no\_data}, cokolwiek powie ekstraktor); cytaty są normalizowane,
        deduplikowane i sprawdzane względem tekstu rozmówcy z tego tematu,
        odrzucane, jeśli to tylko znaczniki zastępcze, czytają się jak polecenie
        do modelu albo ponowna kontrola PII je oflaguje; pewność tematu
        sprzecznego jest ograniczona do \texttt{medium}.
  \item Kanoniczny JSON przechodzi przez \texttt{RecordValidator}. Rekord jest
        \emph{zgłaszalny} (\texttt{Submittable}), gdy wywiad się zakończył,
        rekord jest ważny i ma co najmniej jeden temat \texttt{covered}
        (rekord z samymi \texttt{no\_data} nic nie mówi i nie jest zgłaszalny).
\end{enumerate}

\subsection{Obrona przed wstrzyknięciem i budżety}

Tekst rozmówcy jest niezaufanymi danymi. Trafia do promptu tylko wewnątrz
\texttt{DataBlock}: jeden obiekt JSON na turę w osobnej linii (odpowiedź nie
sfałszuje nagłówka tury ani znacznika), między znacznikami z losowym 64-bitowym
nonce. Test odtwarza odpowiedź fałszującą znacznik końca i linię roli i sprawdza,
że każdy prompt ma dokładnie jeden znacznik początku i końca. Agent nie
udostępnia żadnemu modelowi narzędzi. Dokumentacja wylicza obronę w głąb:
(1) rozdzielenie bloku danych, (2) przepływem steruje maszyna, nie model,
(3) \texttt{QuestionGuard} i zastępcze brzmienie protokołu, (4) wyjście
ekstraktora jest walidowane schematem i przeliczane przez kod, (5) cytaty muszą
istnieć w słowach rozmówcy, a te czytające się jak polecenia są odrzucane,
(6) zamknięty schemat rekordu, (7) zdarzenia \texttt{injection.suspected}
czynią próby widocznymi dla zestawu ocen.

\uwaga{Ryzyka resztkowe nie są ``zaklinane'': prawdziwy ekstraktor nadal może
zostać namówiony do zawyżenia oceny tematu, który ma prawdziwe cytaty
(cytaty ograniczają dowody, nie osąd); listy sygnałów są angielskie
z kilkoma polskimi frazami; wierność cytatu nie oznacza prawdziwości
(rozmówca mógł skłamać).}

Budżety (w protokole, \texttt{limits}): 40 tur prowadzącego, 60 wywołań modelu,
60\,000 szacowanych tokenów, 2000 znaków odpowiedzi, próg lakoniczności
3 słowa i seria 3, 2 wrogie odpowiedzi do zamknięcia, po 1 dopytaniu,
wyjaśnieniu i przekierowaniu na temat, 2 prośby o zgodę. Są to wartości
startowe, a nie zmierzone optimum. Osiągnięcie budżetu zamyka wywiad
łagodnie, a ekstrakcja nadal się odbywa. Z prawdziwym dostawcą nad tym stoi
\emph{twardy} pułap (dwukrotność tokenów) zatrzymujący następne wywołanie.

\subsection{Czego heurystyki potrafią, a czego nie}

\texttt{ReplyAnalyzer} jest oparty na regułach celowo: decyzje chroniące
rozmówcę muszą być testowalne i nie zależeć od ``nastroju'' modelu. Koszt to
topornie czytanie: mglista to odpowiedź o 4--25 słowach z ogólnikiem
(``w porządku'', ``jakoś'') i bez konkretnej wskazówki (cyfra, ``na przykład'',
``ponieważ''); sprzeczność wykrywa się porównując kierunek słów pozytywnych
i negatywnych między kolejnymi odpowiedziami, więc sarkazm i negacja ją mylą.
Nazwiska pochodzą z detektora PII; w trybie \texttt{FailClosed} maskuje on
także nierozpoznane słowa pisane wielką literą, co kosztuje naturalność
rozmowy, a kupuje czułość. Mierzonej dokładności tych reguł jeszcze nie ma;
klasyfikator wspierany modelem to kandydat na dalszą pracę.

\subsection{Persony i skryptowy model testowy}

Osiem symulowanych rozmówców to dane w \texttt{src/}\allowbreak\texttt{ExitInterviewAgent.}\allowbreak\texttt{Personas/Data},
sprawdzane schematem. Wszystkie są syntetyczne (pracodawca to oczywiście
fikcyjny \texttt{widgetron-ltd}, a nazwiska i adresy zmyślone, na zarezerwowanych
domenach). Wybór wariantów odpowiedzi zależy od ziarna (SplitMix64), więc jest
stabilny na każdej platformie.

\begin{tabularx}{\linewidth}{@{}lXX@{}}
\toprule
Persona & Zachowanie & Oczekiwany wynik \\
\midrule
\texttt{talkative} & długie, konkretne odpowiedzi & wszystkie tematy pokryte, zgłaszalny \\
\texttt{terse} & jedno do trzech słów & brak pokrytego tematu, niezgłaszalny \\
\texttt{hostile} & jedna użyteczna odpowiedź, potem wrogość & zamknięcie po jednym przyjęciu; wrogich zdań nie cytuje się \\
\texttt{vague} & ogólniki, potem szczegóły po dopytaniu & sześć dopytań (po jednym na temat) \\
\texttt{names-manager} & podaje nazwisko przełożonego, e-mail i telefon & zamaskowane; jedno przekierowanie; podłożone dane nigdzie nie wyciekają \\
\texttt{prompt-injection} & ładunki wymierzone w prowadzącego, ekstraktor i sędziego & protokół, tematy i kształt rekordu bez zmian \\
\texttt{withdraws-consent} & dwa tematy, potem wycofuje zgodę & brak zapisu rozmowy i rekordu \\
\texttt{contradictory} & opisuje zarządzanie, a potem odwrotnie & jedno wyjaśnienie; pewność ograniczona \\
\bottomrule
\end{tabularx}

\medskip
\texttt{ScriptedChatClient} to model offline: gra prowadzącego i dopytującego,
zwracając tekst zalążkowy protokołu, a ekstraktora --- dzieląc zdania
rozmówcy, wybierając najbardziej konkretne na cytaty i oceniając liczeniem słów
pozytywnych i negatywnych. Nie ma sieci, poświadczeń ani losowości.
\finding{To element testów i demonstracji, a nie wzorzec jakości.} Pokazuje,
że instalacja działa od końca do końca (maszyna stanów, maskowanie, blok
danych, schemat, krok cytatów, walidacja, śledzenie, determinizm i niezmienniki
przy każdej personie), ale \emph{nie} pokazuje, czy prawdziwy model formułuje
neutralne pytania, opiera się wstrzyknięciu, wiernie wyciąga dane ani czy jego
oceny są użyteczne. Każda liczba z przebiegu z modelem testowym opisuje ten
model.

\subsection{Dostawcy modeli}

Moduł \texttt{Providers} udostępnia prawdziwe modele za
\texttt{Microsoft.Extensions.AI.IChatClient}. Dokument
\repopath{docs/architecture/providers.md} (T6): zaimplementowane i
zweryfikowane względem atrap, ale \textbf{nie wykonano żadnego wywołania
prawdziwego dostawcy} (brak klucza i ograniczona sieć), a projekt niczego nie
mierzy o jakości prawdziwych modeli.

\begin{longtable}{@{}p{0.2\linewidth}p{0.4\linewidth}p{0.3\linewidth}@{}}
\toprule
\texttt{--provider} & Backend i pakiet & Klucz \\
\midrule
\endhead
\texttt{anthropic} & API Anthropic z własnym kluczem; oficjalny pakiet
  \texttt{Anthropic} & \texttt{ANTHROPIC\_API\_KEY} (wymagany) \\
\texttt{openai-compatible} (alias \texttt{openai}) & dowolny punkt końcowy
  zgodny z czatem OpenAI (hostowany lub lokalna brama); adapter
  \texttt{Microsoft.Extensions.AI.OpenAI} & \texttt{OPENAI\_API\_KEY}
  (wymagany dla \texttt{api.openai.com}, opcjonalny dla własnej bramy) \\
\texttt{ollama} & lokalny serwer Ollama, natywne \texttt{/api/chat}; klient
  pisany ręcznie (ADR-0032) & brak \\
\texttt{mock} & skryptowy model offline & brak \\
\bottomrule
\end{longtable}

Ustawienia pochodzą kolejno z: flagi, zmiennej środowiskowej, pliku
konfiguracyjnego i wartości domyślnej; dostawca i model nie mają wartości
domyślnych, a identyfikator modelu nigdzie nie jest wpisany na stałe.
Sam klucz pochodzi \emph{wyłącznie} ze zmiennej środowiskowej: nie ma flagi
\texttt{--api-key}, a plik konfiguracyjny z kluczem jest odrzucany. Klucz nigdy
nie jest drukowany, logowany, śledzony ani wstawiany do błędu.

\textbf{Protokół zaufania.} Dostawca widzi całą rozmowę (każde pytanie i każdą
odpowiedź, w wywołaniu ekstraktora naraz), na warunkach \emph{twojego} konta,
których projekt nie weryfikował. CLI mówi o tym i prosi o wpisanie \texttt{yes}
przed pierwszym wywołaniem zewnętrznego dostawcy (\texttt{--yes-i-understand}
dla skryptów); model na własnej maszynie dostaje krótszą informację bez pytania.
Odpowiedzi są maskowane przy odbiorze, więc rozpoznane nazwisko nie dociera do
dostawcy w kolejnych wywołaniach (strażnik jest heurystyczny). Zdalny serwer
Ollama (adres LAN) traktuje się jak zewnętrzny, a ``lokalna'' brama zgodna z
OpenAI może sama przekazywać dalej do zdalnego dostawcy --- tego program nie
widzi. Nic nie jest zapisywane na dysk, o ile o to nie poprosisz
(\texttt{--out}, \texttt{--save-transcript}).

\textbf{Zachowanie przy awarii.} Kody 429, 5xx, przekroczenia czasu i błędy
połączenia są ponawiane (domyślnie do 3 razy z losowanym wycofywaniem,
\texttt{Retry-After} honorowany do 30~s). Kody 401, 403, inne 4xx i
przekierowanie są \emph{fatalne}: wywiad się zatrzymuje, nic nie zostaje,
kod wyjścia 5. Awaria jednego wywołania po ponowieniach oznacza, że ta tura
używa własnego brzmienia protokołu; trzy takie wywołania z rzędu są fatalne.
Ctrl-C lub koniec wejścia: rozmówca wyszedł, zapis odrzucony, nic zapisane.
Cena nie jest wbudowana (nie da się jej zweryfikować stąd): koszt liczy się
tylko, gdy podasz \texttt{--price-in} i \texttt{--price-out}.

\textbf{Czego odmawia i dlaczego.} Poświadczeń subskrypcji Claude (nazwanie
\texttt{CLAUDE\_CODE\_}\allowbreak\texttt{OAUTH\_TOKEN} jako źródła klucza jest odrzucane
z odesłaniem do README --- ograniczenie: rozpoznawalnego formatu takiego tokenu
nie udokumentowano, więc token wklejony do \texttt{ANTHROPIC\_API\_KEY} nie
zostanie rozpoznany i odrzuci go dopiero API) oraz GitHub Copilot (brak
adaptera, a adres bazowy w domenie \texttt{githubcopilot.com} jest odrzucany;
nazwa hosta z ogólnej wiedzy, niezweryfikowana). Telemetria OpenTelemetry jest
domyślnie \emph{wyłączona}; po włączeniu eksportuje tylko źródła agenta
i dostawców, nigdy treści.

Dodanie dostawcy (P10) to implementacja \texttt{IChatClient} i trzy drobne
zmiany: składowa \texttt{Provider}\allowbreak\texttt{Kind}, wiersz \texttt{Provider}\allowbreak\texttt{Catalog} i jeden
\texttt{case} w \texttt{Provider}\allowbreak\texttt{ChatClients.Create}; te same testy obejmują
nowy typ.

\section{Złożenie zgłoszenia i paragon}

\subsection{Zasada: wysyłka jest osobna, opcjonalna i nigdy automatyczna}

Wywiad kończy się plikiem rekordu (\texttt{interview --out}). Złożenie go na
serwerze to osobna decyzja. Wszystkie ścieżki (portal, MCP, CLI) kończą się w
jednej usłudze aplikacyjnej, \texttt{SubmissionService}, więc reguły są
wszędzie te same.

\begin{longtable}{@{}p{0.34\linewidth}p{0.12\linewidth}p{0.2\linewidth}p{0.26\linewidth}@{}}
\toprule
Punkt końcowy & Kto & Autoryzacja & Sekret w \\
\midrule
\endhead
\texttt{POST /api/v1/submissions} & portal (BFF) & polityka \texttt{account} & \texttt{Authorization} \\
narzędzie MCP \texttt{submit\_interview\_record} & host MCP & token MCP, zakres \texttt{interview:submit} & \texttt{Authorization} \\
\texttt{POST /api/v1/tickets} & portal (BFF) & polityka \texttt{account} & \texttt{Authorization} \\
\texttt{POST /api/v1/submissions/ticketed} & CLI & brak (anonimowy) & nagłówek \texttt{X-Submission-Ticket} \\
\texttt{DELETE /api/v1/receipts} & każdy z kodem & brak (anonimowy) & nagłówek \texttt{X-Receipt-Code} \\
\bottomrule
\end{longtable}

\subsection{Bilet: jak CLI zgłasza bez logowania}

Ponieważ CLI nie ma logowania OAuth, konto wystawia mu jednorazowy \emph{bilet}.
W portalu na stronie \texttt{/cli} (po zalogowaniu, ``Create a ticket'') powstaje
losowy, krótkotrwały (15 minut domyślnie; wygaśnięcie zaokrąglane w górę do 5 minut, więc w praktyce 15--20), jednorazowy bilet, \emph{niezwiązany
z żadnym pracodawcą}. Pokazywany jest raz. Serwer trzyma tylko jego skrót.
Aby ograniczyć nadużycia, dozwolone są 3 niewykorzystane bilety na konto i 10
wystawień na godzinę. Realizacja biletu daje serwerowi identyfikator konta
(\texttt{sub}) jednokrotnie, z którego liczy się wpis księgi; potem wiersz
biletu jest usuwany. Nie wymaga to żadnej zmiany w \texttt{authservice}.

\finding{Zastrzeżenie, które CLI pokazuje użytkownikowi przed potwierdzeniem:
chwila realizacji biletu jest punktem korelacji --- bilet wskazuje konto
webowe, więc operator może zobaczyć, że zgłoszenie rekordu nastąpiło razem
z realizacją biletu (T-09 w modelu zagrożeń, zaakceptowane w briefie).}

\begin{figure}[H]
\centering
\includegraphics[width=0.8\linewidth,height=0.55\textheight,keepaspectratio]{\dgm{p1-zgloszenie}}
\caption{Zgłoszenie z CLI i usunięcie przez paragon (uproszczenie diagramów z
\texttt{cli-submission.md} i \texttt{submission-flow.md}).}
\end{figure}

\subsection{Co robi CLI, krok po kroku}

\begin{enumerate}[itemsep=2pt]
  \item \texttt{submit --record plik --server adres}: \textbf{ponowna kontrola
        lokalna} rekordu (schemat, ujawnienie AI, co najmniej jeden pokryty
        temat, skan PII); każda wątpliwość = odmowa.
  \item \textbf{Podgląd dokładnie tego, co opuści maszynę}, wraz z samym
        rekordem. Użytkownik wpisuje słowo \texttt{submit}. Dopiero wtedy
        program prosi o bilet: ukryty monit, zmienna
        \texttt{EXIT\_}\allowbreak\texttt{INTERVIEW\_}\allowbreak\texttt{TICKET} albo jedna linia standardowego wejścia.
        \textbf{Nie ma flagi z biletem ani z kodem paragonu} (argv trafia do
        listy procesów i historii powłoki); pilnuje tego test architektury.
  \item Jedno żądanie \texttt{POST} na \texttt{/api/v1/submissions/ticketed}
        z rekordem jako ciałem i biletem w nagłówku.
\end{enumerate}

Z maszyny wychodzi: rekord bajt w bajt tak, jak pokazano, bilet w jednym
nagłówku oraz to, co wysyła każdy klient HTTP (adres klienta, czas,
\texttt{User-Agent: exit-interview} bez wersji, ścieżka). Nie wychodzi: zapis
rozmowy, prompty i odpowiedzi, nazwy plików, dostawca, model ani klucz API,
ciasteczka, \texttt{Authorization}, zapytanie w adresie.

Adres serwera pochodzi z \texttt{--server} lub
\texttt{EXIT\_}\allowbreak\texttt{INTERVIEW\_}\allowbreak\texttt{SERVER\_URL}, \textbf{bez wartości domyślnej} (nic nie
jest wdrożone); wymagany jest \texttt{https} (\texttt{http} tylko dla
loopback); przekierowania nigdy nie są podążane; po wysłaniu choć jednego
bajtu nic się nie ponawia; 10~s na połączenie i 30~s łącznie; odpowiedzi
ponad 64~KiB są odrzucane.

\subsection{Co serwer robi z rekordem}

Serwer wykonuje kolejno: odczyt ograniczony (limit 160~KiB) i limity szybkości;
dla ścieżki z biletem --- odczyt skrótu biletu tylko do odczytu (daje \texttt{sub},
niczego nie zużywa); walidację rozmiaru, schematu, zduplikowanych kluczy
i \texttt{piiMasked}; \emph{ponowny skan PII} każdego cytatu i pola
tekstowego (każdy wyjątek lub przekroczenie czasu = odrzucenie); wymóg
\texttt{aiDisclosed = true}; zapytanie weryfikatora zatrudnienia
(\emph{nigdy nie dostaje rekordu}; odpowiedź ``niedostępny'' obniża poziom do
\texttt{Unchecked} i zgłoszenie trwa dalej); wreszcie jedną transakcję: sprawdza
wolność identyfikatora wywiadu i znacznika księgi, usuwa bilet (jeden
usunięty wiersz = ten wywołujący wygrał), wstawia wpis księgi, rekord
i paragon. Utrata wyścigu o unikalny indeks wycofuje wszystko.

\subsection{Tabele i to, z czym da się je (nie) połączyć}

\begin{longtable}{@{}p{0.2\linewidth}p{0.4\linewidth}p{0.33\linewidth}@{}}
\toprule
Tabela & Kolumny (skrót) & Cel \\
\midrule
\endhead
\texttt{Records} & identyfikator wywiadu, \texttt{EmployerRef}, kanoniczny JSON,
  poziom weryfikacji (\texttt{Verified}/\texttt{Unverified}/\texttt{Unchecked}),
  \texttt{CreatedWeek} (poniedziałek) & sam rekord; tydzień służy tylko do
  usuwania po wieku \\
\texttt{SubmissionLedger} & losowe id, \texttt{KeyId}, \texttt{Tag} (HMAC,
  unikalny), \texttt{CreatedWeek} & jeden wpis na pracodawcę na konto \\
\texttt{Receipts} & losowe id, \texttt{CodeHash} (SHA-256, unikalny),
  \texttt{RecordId} (klucz obcy, kaskada) & usuwanie przez kod paragonu \\
\texttt{SubmissionTickets} & losowe id, \texttt{TokenHash}, \texttt{Sub},
  \texttt{ExpiresAt} & bilety CLI; usuwane przy realizacji lub wygaśnięciu \\
\bottomrule
\end{longtable}

Niezmienniki (każdy ma test): \emph{żadna tabela poza biletem nie trzyma
\texttt{sub}; księga nie ma kolumny ani klucza obcego odnoszącego się do
rekordu; żadna tabela nie trzyma jednocześnie wartości pochodzącej z konta
i wartości identyfikującej rekord; każdy klucz jest losowy; jedyne
przechowywane znaczniki czasu to pasma tygodnia (plus wygaśnięcie biletu).}
Księga trzyma HMAC pary (konto, pracodawca) pod kluczem, którego baza nie ma;
klucz jest rotowalny, a znacznik sprawdza się dla każdego aktywnego klucza.

\uwaga{Czego \emph{nie} ukryto. (1) Wiersze zapisane w jednej transakcji mają
wspólny identyfikator transakcji PostgreSQL i leżą blisko siebie w stercie:
kto ma bezpośredni dostęp do plików lub kopii fizycznej, może sparować wpis
księgi z rekordem, choć żadna kolumna ich nie łączy (T-08, T-09, OP-13).
(2) Operator z kluczem HMAC i bazą może obliczyć znacznik dla dowolnego
konta i pracodawcy i potwierdzić zgłoszenie. (3) Kto widzi żywy ruch, widzi
jednoczesną realizację biletu i napływ rekordu. (4) Usunięcie rekordu przez
paragon nie czyści wpisu księgi (nic ich nie łączy), więc konto pozostaje
zamknięte dla tego pracodawcy do końca okna księgi (365 dni).}

\subsection{Paragon i usuwanie}

Po udanym zgłoszeniu serwer zwraca \emph{kod paragonu} --- losowy, pokazywany
\textbf{raz}. Serwer trzyma tylko jego skrót SHA-256. Okazanie kodu przez
\texttt{DELETE /api/v1/receipts} usuwa rekord bez łączenia go z kontem;
odpowiedź \texttt{204} jest taka sama dla każdego poprawnie sformowanego kodu,
więc nie potwierdza, że rekord istniał (stąd komunikat CLI: ``jeśli rekord
o tym kodzie istniał, został właśnie usunięty''). Źle sformowany kod to
\texttt{400 INVALID\_RECEIPT\_CODE} (literówka, a nie usunięcie). Odpowiedź ma
dolny próg czasu (150~ms), a punkt końcowy limity szybkości (6 na IP na minutę,
60 globalnie na minutę).

CLI przyjmuje kod ze zmiennej \texttt{EXIT\_}\allowbreak\texttt{INTERVIEW\_}\allowbreak\texttt{RECEIPT\_CODE},
ukrytego monitu lub standardowego wejścia --- nigdy z argumentu. Świeżo wydany
kod pojawia się na standardowym wyjściu raz, plus opcjonalnie w pliku
\texttt{--save-receipt} (tryb 0600, nigdy nie nadpisuje istniejącego pliku).
\finding{Zgubiony kod nie jest odzyskiwalny: nikt nie może go dla ciebie
wyszukać}, bo rekord nie ma powiązania z kontem, a usunięcie konta nie sięga
rekordów (OP-12). W trybie MCP kod przechodzi przez hosta, więc Claude i jego
dostawca mogą go przeczytać; każdy, kto go ma, może usunąć rekord.

\subsection{Wyniki i kody odrzuceń}

Odpowiedzi błędów to RFC 9457 (\texttt{application/problem+json}); żaden kod ani
komunikat nie niesie przesłanego tekstu. Najważniejsze:

\begin{longtable}{@{}p{0.1\linewidth}p{0.3\linewidth}p{0.55\linewidth}@{}}
\toprule
HTTP & Kod & Znaczenie \\
\midrule
\endhead
201 & (ciało: kod paragonu) & zapisano; kod nigdy więcej nie jest pokazywany \\
401 & \texttt{TICKET\_INVALID} & nieznany, wygasły, zużyty, źle sformowany albo brak biletu --- jedna odpowiedź \\
403 & \texttt{EMPLOYMENT\_NOT\_VERIFIED} & tylko przy włączonej polityce ścisłej \\
409 & \texttt{ALREADY\_SUBMITTED} & to konto już zgłosiło dla tego pracodawcy w oknie \\
409 & \texttt{INTERVIEW\_ID\_TAKEN} & rekord jest już zapisany (zwykle wcześniejsza próba, której odpowiedź zginęła) \\
413 & \texttt{PAYLOAD\_TOO\_LARGE} & ponad 160~KiB \\
422 & kody biblioteki rekordu & naruszenie schematu, z listą (kod, ścieżka) \\
422 & \texttt{AI\_NOT\_DISCLOSED} & nie powiedziano rozmówcy, że prowadzący jest AI \\
422 & \texttt{PII\_DETECTED} & ponowny skan znalazł dane osobowe (tylko rodzaje, nigdy treść) \\
422 & \texttt{PII\_CHECK\_FAILED} & skan nie mógł się zakończyć; nic nie zapisano \\
429 & \texttt{TICKET\_LIMIT} / \texttt{rate\_limited} & zbyt wiele żywych biletów albo limit szybkości \\
\bottomrule
\end{longtable}

Kody wyjścia CLI: 0 złożono, 2 błąd użycia, 3 niepotwierdzone lub brak biletu,
4 rekord nie przeszedł lokalnej kontroli, 5 serwer odmówił, 6 błąd sieci lub
nieznany wynik, 130 anulowano. Wynik jest \emph{nieznany} przy przekroczeniu
czasu, zerwanym połączeniu lub anulowaniu po wysłaniu bajtu: serwer mógł
zapisać rekord, a paragon przepadł (OP-24). Dlatego CLI niczego nie ponawia po
wysłaniu bajtu.

\subsection{Retencja}

Usługa tła (\texttt{RetentionService}) co 5 minut (domyślnie) usuwa partiami po
500: rekordy (wraz z paragonami) starsze niż 24 miesiące (założenie, ADR-0019),
wpisy księgi starsze niż 365 dni (ADR-0028) i wygasłe bilety. Wiek liczy się do
końca tygodniowego pasma. Loguje wyłącznie liczby.

\subsection{Tryb A: złożenie przez narzędzie MCP}

Serwer MCP (\texttt{interview-service}, ścieżka z \texttt{Mcp:Resource}, np.\
\texttt{/mcp}; oficjalny pakiet \texttt{ModelContext}\allowbreak\texttt{Protocol.AspNetCore}, Streamable
HTTP, bezstanowy) publikuje kontrakt przypięty migawką testową:

\begin{longtable}{@{}p{0.14\linewidth}p{0.36\linewidth}p{0.45\linewidth}@{}}
\toprule
Rodzaj & Nazwa & Uwagi \\
\midrule
\endhead
Prompt & \texttt{conduct\_exit\_}\allowbreak\texttt{interview(language?,} \texttt{employerHint?)} & jedna
  wiadomość: zasady zaufania, dosłowne otwarcie, sześć tematów, reguły pytań
  i rekordu, kolejność: waliduj, potwierdź, złóż \\
Narzędzie & \texttt{validate\_}\allowbreak\texttt{interview\_record} & tylko do odczytu, idempotentne,
  próba na sucho; kody, ścieżki, rodzaje PII \\
Narzędzie & \texttt{submit\_}\allowbreak\texttt{interview\_record} & nie tylko do odczytu, nie
  idempotentne; kod paragonu raz \\
Zasoby & \texttt{exit-interview://}\allowbreak\texttt{schema/}\allowbreak\texttt{record/v1},
  \texttt{.../protocol/v1}, \texttt{.../topics/v1} & schemat, protokół, tematy \\
\bottomrule
\end{longtable}

{\sloppy Użytkownik łączy Claude (własny konektor z adresem, identyfikatorem klienta
i sekretem), a Claude uwierzytelnia się w \texttt{authservice}
(kod autoryzacji z PKCE; token z odbiorcą \texttt{/mcp} i zakresem
\texttt{interview:submit}). Rozmowę prowadzi model-gospodarz,
według promptu; serwer widzi \texttt{sub}, \texttt{client\_id}, zakres, argumenty promptu
(język, wskazówka o pracodawcy) i rekord (najwyżej dwukrotnie: walidacja
i zgłoszenie), ale nie zapis rozmowy ani adres e-mail. Ochrony serwera: dwa
schematy JWT, zakres sprawdzany na poziomie polityki HTTP i ponownie dla
każdego wywołania narzędzia, lista dozwolonych \texttt{Origin}, kontrola wersji
protokołu i rozmiaru, brak sesji i starszego SSE, zamknięte słownictwo dla nazw
wybieranych przez klienta, podłoga logowania SDK, limit szybkości na konto, a
także test, że ścieżka nie używa próbkowania (\emph{sampling}), elicytacji ani
katalogów \emph{roots}.\par}

\finding{Granica zaufania trybu A jest także jego słabością: wywiad prowadzi
model-gospodarz, nie ten projekt.} Czy ujawnia, że jest AI, uzyskuje jasne ``tak'',
przerywa na żądanie, zadaje neutralne pytania i używa dosłownych cytatów, zależy
od modelu wykonującego instrukcje, a serwer niczego z tego nie obserwuje.
Sprawdza tylko to, co może: schemat, zamknięte pola, skan PII, flagę
\texttt{aiDisclosed} (nie wie, czy to faktycznie powiedziano), księgę. Nie
sprawdza nic w kwestii dosłowności cytatów (nie ma zapisu), zgody, przerwania,
neutralności i dopytań. Dlatego powstał zestaw ocen: te same persony mają być
uruchamiane przez hosty trybu A i oceniane tymi samymi miarami; \textbf{dziś ta
liczba nie istnieje --- wierność jakiegokolwiek hosta jest niezmierzona}.
Nie uruchomiono żadnego prawdziwego klienta Claude, a przepływ konektora,
wykrywalność promptu i zachowanie hosta są niezweryfikowane (OP-18).

\subsection{Podsumowanie części}

Produkt składa się z kilku świadomych decyzji: model przynosi użytkownik;
program, a nie model, steruje wywiadem; rekord jest mały, zamknięty i
pozbawiony identyfikatorów; serwer widzi rekord, nie rozmowę; księga i paragon
są od siebie odcięte; a wszystko, czego nie zmierzono (żywi dostawcy, żywy
klient Claude, prawdziwe TLS, prawdziwy ruch, prawdziwa weryfikacja
zatrudnienia), jest nazwane. Kolejne części przewodnika pokazują, jak z tego
korzystać, co dokładnie chroni dane i jak mierzy się jakość wywiadu.
